\uuid{H1ak}
\titre{Module/argument}
\auteur{bodin ; hanani ; mzari}
\organisation{exo7}
\licence{CC BY-NC-SA 4.0}
\datecreate{2020-01-20}
\niveau{L1}
\module{Algèbre}
\chapitre{Nombres complexes}
\sousChapitre{Forme cartésienne, forme polaire}
\difficulte{3}
\theme{Écritures algébrique et géométrique, module, argument}
\exo7qcmid{lille-1.nombres-complexes.014}

\contenu{
\begin{question}
\qclassification{104.01}
\qtags{motcle=module/argument}
Soit $z=e^{i\theta} + e^{i\phi} ,\theta, \phi \in \Rr$ tels que $-\pi < \theta - \phi < \pi$. Quelles sont les assertions vraies ?
\begin{answers}
    \bad{$|z|=2 $}

    \good{$|z|=2\cos (\frac{ \theta -\phi}{2}) $}

    \bad{$\arg z = \theta +\phi  \, [2\pi]$}

    \good{$\arg z = \frac{\theta+ \phi}{2} \, [2\pi]$}

\end{answers}
\begin{explanations}
$ z=e^{i\frac{\theta+\phi}{2}} (e^{i\frac{\theta-\phi}{2}} + e^{i\frac{\phi - \theta}{2}}) = 2 \cos (\frac{\theta-\phi}{2}) e^{i\frac{\theta+\phi}{2}}$. Comme $\theta-\phi \in ]-\pi,\pi[$,   $\cos (\frac{\theta-\phi}{2})>0$. On déduit que : $|z|=2\cos (\frac{\theta-\phi}{2})$ et $\arg z =  \frac{\theta+\phi}{2} \, [2\pi]$.
\end{explanations}
\end{question}
}
